\documentclass{situationist-preprint}

\title{Positive Illusions}
\msid{EXI-18-b}
\msfamily{EXI-18}
\msdiscipline{Meaning, Death and the Absurd}
\mstype{Consequence}
\mskeywords{positive illusions; self-enhancement; optimism; illusion of control; well-being}

\begin{document}
\maketitle

\begin{abstract}
Taylor and Brown proposed that unrealistically positive self-evaluations, exaggerated perceptions of control and unrealistic optimism are features of normal mental health. We derive from the results of EXI-18-a what follows for these three illusions: some of them may be psychologically necessary, none is made true or justified by that necessity, and none may be cultivated beyond what one cannot avoid.
\end{abstract}

\section{The consequence}

\citet{Taylor1988} argued that three positive illusions are associated with persistence, contentment and care for others, and so with mental health; \citet{Colvin1994} disputed the evidence and the conception of health it assumes. EXI-18-a need not settle the dispute.

\begin{proposition}\label{p:pos}
(i) Whether a given positive illusion is psychologically necessary for a given agent is an empirical question, and the answer may differ across agents and across the three kinds. (ii) Even where one is necessary, it is not thereby true or justified. (iii) Its necessity licenses neither adopting it nor cultivating it further, for instance by therapies or practices aimed at inflating self-assessment beyond the evidence.
\end{proposition}

\begin{proof}
(i) follows from Definition~2.1 and Proposition~2.2 of EXI-18-a, since necessity is relative to an agent's capacity to function and the second conjunct of Proposition~2.2 says that not every illusion is necessary. (ii) is Theorem~2.3 of EXI-18-a applied to the three kinds. (iii) is Corollary~2.4 of EXI-18-a: cultivation through selective attention produces belief on insufficient evidence, and each illusion is a respect in which one's life is out of contact with how things are.
\end{proof}

Part (iii) tells against interventions that aim past the evidence, not against those that correct unduly negative self-assessment. And by Remark~2.5 of EXI-18-a, the functions of optimism and of the illusion of control can be served by hope proportioned to probability (EXI-15) and by practical commitment, which do not misrepresent; to that extent the range of necessary positive illusions shrinks.

\begin{thebibliography}{9}\small
\bibitem[Colvin and Block(1994)]{Colvin1994} Colvin, C.~R. and Block, J. (1994). Do positive illusions foster mental health? An examination of the Taylor and Brown formulation. \emph{Psychological Bulletin}, 116(1), 3--20.
\bibitem[Taylor and Brown(1988)]{Taylor1988} Taylor, S.~E. and Brown, J.~D. (1988). Illusion and well-being: A social psychological perspective on mental health. \emph{Psychological Bulletin}, 103(2), 193--210.
\end{thebibliography}

\end{document}
